\begin{table}[!htbp]
\caption{Transfer results under different adaptation ratios.}
\label{tab:4-9}
\centering
\zihao{-5}
\setlength{\tabcolsep}{2.4pt}
\renewcommand{\arraystretch}{1.04}
\begin{ThreePartTable}
\begin{tabularx}{0.90\linewidth}{>{\raggedright\arraybackslash}p{0.25\linewidth}>{\centering\arraybackslash}p{0.13\linewidth}>{\centering\arraybackslash}p{0.10\linewidth}>{\centering\arraybackslash}p{0.10\linewidth}>{\centering\arraybackslash}p{0.10\linewidth}>{\centering\arraybackslash}p{0.10\linewidth}}
\toprule
\makecell[l]{Source domain\\(training cell)} & Metric & 10\% & 30\% & 50\% & 70\% \\
\midrule
\makecell[l]{CALCE CS2\\(CS2\_36)} & $R^2$      & -1.3934 & -0.2679 & 0.5094 & 0.7650 \\
                      & MAE        & 0.0831  & 0.0590  & 0.0366 & 0.0252 \\
                      & MAPE       & 0.1015  & 0.0722  & 0.0445 & 0.0305 \\
                      & RMSE       & 0.0942  & 0.0685  & 0.0426 & 0.0293 \\
\addlinespace[1pt]
\makecell[l]{CALCE CX2\\(CX2\_36)} & $R^2$      & -3.2217 & -1.6085 & -0.3836 & -0.1268 \\
                      & MAE        & 0.1095  & 0.0835  & 0.0614  & 0.0559 \\
                      & MAPE       & 0.1339  & 0.1024  & 0.0747  & 0.0677 \\
                      & RMSE       & 0.1251  & 0.0982  & 0.0716  & 0.0646 \\
\bottomrule
\end{tabularx}
\begin{tablenotes}[flushleft]
\footnotesize
\item[] \parbox[t]{0.98\linewidth}{Note: Results are averaged over three independent runs.}
\end{tablenotes}
\end{ThreePartTable}
\end{table}
